\emph{(i).} By \claimref{std-algebra.quaternions} the elements $1, B_{1}, B_{2}, B_{3}$ multiply as the quaternion units: $B_{k}^{2} = -1$, $B_{1}B_{2} = B_{3} = -B_{2}B_{1}$ and cyclically. The reverse fixes $1$ and sends $B_{k}$ to $-B_{k}$, so $\tilde{x} = x_{0} - x_{1}B_{1} - x_{2}B_{2} - x_{3}B_{3}$ and the cross terms of $x\tilde{x}$ cancel in pairs: $x\tilde{x} = x_{0}^{2} + x_{1}^{2} + x_{2}^{2} + x_{3}^{2}$. The reverse is an anti-automorphism and $N(y)$ is a scalar, so $N(xy) = xy\tilde{y}\tilde{x} = N(y)\,x\tilde{x} = N(x)N(y)$. A sum of four rational squares vanishes only if each does, so $N(x) = 0$ forces $x = 0$, and $x\,\tilde{x}/N(x) = 1$ otherwise. An element with an inverse is no zero divisor; \claimref{synthesis.division-gives-the-rationals} records the same for $\Cl_{\Q}(0,2)$, whose generators $f_{1}, f_{2}$ and their product multiply as $B_{1}, B_{2}, B_{3}$.

\emph{(ii).} Let $S$ be the set of positive integers of the form $N(x)$ with $x$ integral. The product of two integral even elements is integral, so by (i) $S$ is closed under products, and $1 = N(1)$, $2 = N(1 + B_{1})$ lie in $S$. It suffices to show that every odd prime $p$ lies in $S$.

\emph{Step 1: some multiple $mp$ with $0 < m < p$ lies in $S$.} The $(p+1)/2$ integers $a^{2}$, $0 \le a \le (p-1)/2$, are pairwise incongruent modulo $p$ (if $a^{2} \equiv a'^{2}$ then $p$ divides $(a - a')(a + a')$, and both factors lie strictly between $-p$ and $p$, with $a + a' = 0$ only for $a = a' = 0$), and so are the $(p+1)/2$ integers $-1 - b^{2}$. Together they are $p + 1$ integers, so two are congruent, one from each family: $a^{2} + b^{2} + 1 = mp$ with $0 < mp \le 2\bigl(\tfrac{p-1}{2}\bigr)^{2} + 1 < p^{2}$, so $0 < m < p$, and $mp = N(a + bB_{1} + B_{2})$.

\emph{Step 2: the least such $m$ is $1$.} Let $m$ be the least positive integer with $mp = N(x)$, $x$ integral; by Step 1, $m < p$. Suppose $m$ is even. Then $x_{0}^{2} + \dots + x_{3}^{2}$ is even, so the number of odd components is $0$, $2$ or $4$, and the components can be put in two pairs of equal parity, say $(x_{0}, x_{1})$ and $(x_{2}, x_{3})$ after renaming. The element $y$ with components $\tfrac12(x_{0} + x_{1})$, $\tfrac12(x_{0} - x_{1})$, $\tfrac12(x_{2} + x_{3})$, $\tfrac12(x_{2} - x_{3})$ is integral and $N(y) = \tfrac12 N(x) = \tfrac{m}{2}p$, against the choice of $m$. So $m$ is odd. Suppose $m > 1$. Choose integers $y_{i} \equiv x_{i} \pmod m$ with $|y_{i}| < m/2$, and let $y$ be the even element with these components. Then $N(y) \equiv N(x) \equiv 0 \pmod m$ and $N(y) < 4(m/2)^{2} = m^{2}$, so $N(y) = mm'$ with $0 \le m' < m$. If $m' = 0$ then $y = 0$, so $m$ divides every $x_{i}$, $m^{2}$ divides $N(x) = mp$, and $m$ divides the prime $p$ with $1 < m < p$: impossible. So $0 < m' < m$. Put $z = x\tilde{y}$, an integral even element. Its components are polynomials with integer coefficients in the $x_{i}$ and $y_{i}$, and $y_{i} \equiv x_{i}$, so they are congruent modulo $m$ to the components of $x\tilde{x} = N(x) = mp$, which are $mp, 0, 0, 0$: every component of $z$ is divisible by $m$. So $z = mz'$ with $z'$ integral, and $m^{2}N(z') = N(z) = N(x)N(y) = m^{2}m'p$, that is $m'p = N(z')$ with $0 < m' < m$, against the choice of $m$. Hence $m = 1$ and $p \in S$.

Every positive integer is a product of primes, so it lies in $S$. For a positive rational $t = a/b$ with $a, b$ positive integers, $ab = N(y)$ with $y$ integral and $N(y/b) = ab/b^{2} = t$.

\emph{(iii).} By (i) every element of $\mathsf{P}$ is a sum of four rational squares, not all zero, hence a positive rational, and by (ii) every positive rational lies in $\mathsf{P}$. The positive rationals are closed under sums, products and inverses and $\Q$ is the disjoint union of them, $0$ and their negatives. Let $\mathsf{P}'$ be an ordering. For $s \ne 0$ either $s$ or $-s$ lies in $\mathsf{P}'$, and in both cases $s^{2} = (\pm s)^{2} \in \mathsf{P}'$; so $\mathsf{P}'$ contains every nonzero square and, being closed under sums, every sum of four squares not all zero, that is $\mathsf{P}$. If some $t \in \mathsf{P}'$ were not in $\mathsf{P}$ then $-t \in \mathsf{P} \subset \mathsf{P}'$, against the disjointness of $\mathsf{P}'$ and $-\mathsf{P}'$. So $\mathsf{P}' = \mathsf{P}$. The last sentence is the definition of the order by its positive cone.

\emph{(iv).} Suppose $a^{2} + b^{2} + c^{2} = 7d^{2}$ with $d \ne 0$ and $|d|$ least. A square is $0$, $1$ or $4$ modulo $8$. If $d$ is odd, $7d^{2} \equiv 7 \pmod 8$, and no three of $0, 1, 4$ sum to $7$ modulo $8$. If $d$ is even, $a^{2} + b^{2} + c^{2} \equiv 0 \pmod 4$; squares are $0$ or $1$ modulo $4$, so $a, b, c$ are all even and $(a/2, b/2, c/2, d/2)$ is a solution with smaller $|d|$. Both cases are impossible. A bivector $B = b_{1}B_{1} + b_{2}B_{2} + b_{3}B_{3}$ has $B^{2} = -N(B) = -(b_{1}^{2} + b_{2}^{2} + b_{3}^{2})$ by (i), so $B^{2} = -7$ with rational $b_{k}$ would give, after clearing a common denominator $d$, a solution of the equation.

\emph{(v).} A rational quadratic form is diagonalised by a rational change of variables (completing squares; when a diagonal entry vanishes and the form is non-negative, the whole row vanishes, since otherwise $r^{T}Ar$ takes a negative value on a suitable combination of two basis vectors). So $A = LDL^{T}$ with $L \in \Mat_{n}(\Q)$ invertible and $D$ diagonal, and $D_{jj} = e_{j}'^{T}A\,e_{j}' \ge 0$ for the vectors $e_{j}' = (L^{T})^{-1}e_{j}$. By (ii) each $D_{jj} = c_{j1}^{2} + \dots + c_{j4}^{2}$ with rational $c_{ji}$, so with $L_{j}$ the $j$-th column of $L$,
\[
A = \sum_{j}D_{jj}L_{j}L_{j}^{T} = \sum_{j}\sum_{i=1}^{4}(c_{ji}L_{j})(c_{ji}L_{j})^{T},
\]
at most $4n$ terms. Conversely $r^{T}yy^{T}r = (y \cdot r)^{2} \ge 0$. In the matrix algebra with the transpose as adjoint, $xx^{\dagger} = xx^{T}$ is the sum of the $y y^{T}$ over the columns $y$ of $x$, and $yy^{T} = xx^{T}$ for the matrix $x$ with first column $y$ and zeros elsewhere; so the finite sums of elements $xx^{\dagger}$ are the finite sums of matrices $yy^{T}$, and these are the non-negative symmetric elements.
